\documentclass[10pt,a4paper]{amsart}
\usepackage{amsmath,amssymb,amsthm,graphicx}
% For commutative diagrams we use Xy-pic
\usepackage[all]{xy}

% We use 2cell for 2-commutative diagrams.
\xyoption{2cell}
% Theorem environments.
%
\theoremstyle{plain}
\newtheorem{theorem}[subsection]{Theorem}
\newtheorem{proposition}[subsection]{Proposition}
\newtheorem{lemma}[subsection]{Lemma}

\theoremstyle{definition}
\newtheorem{definition}[subsection]{Definition}
\newtheorem{example}[subsection]{Example}
\newtheorem{exercise}[subsection]{Exercise}
\newtheorem{situation}[subsection]{Situation}

\theoremstyle{remark}
\newtheorem{remark}[subsection]{Remark}
\newtheorem{remarks}[subsection]{Remarks}

\numberwithin{equation}{subsection}

% Macros
%
\def\lim{\mathop{\mathrm{lim}}\nolimits}
\def\colim{\mathop{\mathrm{colim}}\nolimits}
\def\Spec{\mathop{\mathrm{Spec}}}
\def\Hom{\mathop{\mathrm{Hom}}\nolimits}
\def\Ext{\mathop{\mathrm{Ext}}\nolimits}
\def\SheafHom{\mathop{\mathcal{H}\!\mathit{om}}\nolimits}
\def\SheafExt{\mathop{\mathcal{E}\!\mathit{xt}}\nolimits}
\def\Sch{\mathit{Sch}}
\def\Mor{\mathop{\mathrm{Mor}}\nolimits}
\def\Ob{\mathop{\mathrm{Ob}}\nolimits}
\def\Sh{\mathop{\mathit{Sh}}\nolimits}
\def\NL{\mathop{N\!L}\nolimits}
\def\CH{\mathop{\mathrm{CH}}\nolimits}
\def\proetale{{pro\text{-}\acute{e}tale}}
\def\etale{{\acute{e}tale}}
\def\QCoh{\mathit{QCoh}}
\def\Ker{\mathop{\mathrm{Ker}}}
\def\Im{\mathop{\mathrm{Im}}}
\def\Coker{\mathop{\mathrm{Coker}}}
\def\Coim{\mathop{\mathrm{Coim}}}

% Boxtimes
%
\DeclareMathSymbol{\boxtimes}{\mathbin}{AMSa}{"02}

%
% Macros for moduli stacks/spaces
%
\def\QCohstack{\mathcal{QC}\!\mathit{oh}}
\def\Cohstack{\mathcal{C}\!\mathit{oh}}
\def\Spacesstack{\mathcal{S}\!\mathit{paces}}
\def\Quotfunctor{\mathrm{Quot}}
\def\Hilbfunctor{\mathrm{Hilb}}
\def\Curvesstack{\mathcal{C}\!\mathit{urves}}
\def\Polarizedstack{\mathcal{P}\!\mathit{olarized}}
\def\Complexesstack{\mathcal{C}\!\mathit{omplexes}}
% \Pic is the operator that assigns to X its picard group, usage \Pic(X)
% \Picardstack_{X/B} denotes the Picard stack of X over B
% \Picardfunctor_{X/B} denotes the Picard functor of X over B
\def\Pic{\mathop{\mathrm{Pic}}\nolimits}
\def\Picardstack{\mathcal{P}\!\mathit{ic}}
\def\Picardfunctor{\mathrm{Pic}}
\def\Deformationcategory{\mathcal{D}\!\mathit{ef}}
\usepackage{cleveref}
\crefname{originaldefinition}{original definition}{original definitions}
\Crefname{originaldefinition}{Original definition}{Original definitions}
\makeatletter
\newlabel{situation-setup}{{02QL}{source}{}{sourcepointer-1}{}}
\newlabel{situation-setup@cref}{{[originaldefinition][0][]02QL}{[1][0][]source}{}{}{}}
\makeatother
\crefname{originallemma}{original lemma}{original lemmas}
\Crefname{originallemma}{Original lemma}{Original lemmas}
\makeatletter
\newlabel{section-pre-derived}{{0ESY}{source}{}{sourcepointer-2}{}}
\newlabel{section-pre-derived@cref}{{[originallemma][0][]0ESY}{[1][0][]source}{}{}{}}
\makeatother
\crefname{originallemma}{original lemma}{original lemmas}
\Crefname{originallemma}{Original lemma}{Original lemmas}
\makeatletter
\newlabel{lemma-exact-sequence-closed-chow}{{0F94}{source}{}{sourcepointer-3}{}}
\newlabel{lemma-exact-sequence-closed-chow@cref}{{[originallemma][0][]0F94}{[1][0][]source}{}{}{}}
\makeatother
\title{Existence of a bivariant class on a closed subscheme characterised by vanishing on the complementary component}
\author{The Stacks Project Authors}
\date{stacks-project:chow.tex:local-snapshot-20261012; GNU FDL 1.2 or later}
\begin{document}
\maketitle
\noindent\textit{Editorial note.} All delivered mathematical text is the author's own, copied from the pinned Stacks Project chapter file chow.tex; no statement or proof step was rewritten, reordered or removed.\par\medskip
\noindent\textit{Editorial note.} Preamble provenance: the pinned chapter file opens with an input of the shared preamble and carries no preamble text of its own, so the project's single Stacks-wide preamble staged at sources/stacks-chow/preamble.tex is excerpted verbatim as two spans, lines 14-18 (Xy-pic setup) and lines 154-217 (the native theorem declarations and mathematical macros). Nothing was added, removed or redefined; the wrapper class and the amsmath/amssymb/amsthm/graphicx base packages are supplied by the assembler.\par\medskip
\noindent\textit{Editorial note.} The counted statement and proof contain no cross reference, no citation and no included file, so the delivered body is self contained.\par\medskip
\noindent\textit{Editorial note.} Canonical location of the counted proposition: Stacks tag 0F9H (sources/stacks-chow/pinned-original-source.tex line 9695).\par\medskip
\noindent\textit{Allowed standard prerequisites (editorial).}\begin{itemize}
\item Writing a cycle class as a sum of pushforwards from a cover by two closed subschemes. Same chapter, Lemma labelled lemma-exact-sequence-closed-chow (source line 2991), invoked at proof lines 9734-9736. Supplies the decomposition alpha' = i'\_\{1,*\}alpha'\_1 + i'\_\{2,*\}alpha'\_2 and the description of the ambiguity of alpha'\_2 as the image of the pushforward from X'\_1 cap X'\_2, which is exactly the ambiguity the hypotheses remove.
\item Chern classes of a perfect object and the capping operation of a bivariant class. Same chapter, Section labelled section-pre-derived (source line 9042), used throughout the proof, and Situation labelled situation-setup (source line 1372) in the statement. Fixes the objects P\_p(E\_2), c\_p(E\_2) and the bivariant Chow groups in which the characterised class lives.
\end{itemize}\medskip
\ifdefined\hypertarget\hypertarget{sourcepointer-1}{}\fi
\noindent\textit{Source reference.} Original definition 02QL: sources/stacks-chow/pinned-original-source.tex, source lines 1372--1372. Same-chapter pointer into the author environment situation at source line 1372; invoked at proof line None of the counted proof. It is retained only as a navigation pointer; every load-bearing step is inside the delivered span.\par\medskip
\ifdefined\hypertarget\hypertarget{sourcepointer-2}{}\fi
\noindent\textit{Source reference.} Original lemma 0ESY: sources/stacks-chow/pinned-original-source.tex, source lines 9042--9042. Same-chapter pointer into the author environment lemma at source line 9042; invoked at proof line 9727 of the counted proof. It is retained only as a navigation pointer; every load-bearing step is inside the delivered span.\par\medskip
\ifdefined\hypertarget\hypertarget{sourcepointer-3}{}\fi
\noindent\textit{Source reference.} Original lemma 0F94: sources/stacks-chow/pinned-original-source.tex, source lines 2991--2991. Same-chapter pointer into the author environment lemma at source line 2991; invoked at proof line 9734 of the counted proof. It is retained only as a navigation pointer; every load-bearing step is inside the delivered span.\par\medskip
\begin{lemma}
\label{lemma-silly}
Let $(S, \delta)$ be as in Situation \ref{situation-setup}. Let $X$ be
locally of finite type over $S$. Let $i_j : X_j \to X$, $j = 1, 2$
be closed immersions such that $X = X_1 \cup X_2$ set theoretically. Let
$E_2 \in D(\mathcal{O}_{X_2})$ be a perfect object. Assume
\begin{enumerate}
\item Chern classes of $E_2$ are defined,
\item the restriction $E_2|_{X_1 \cap X_2}$ is zero,
resp.\ isomorphic to a finite locally free $\mathcal{O}_{X_1 \cap X_2}$-module
of rank $< p$ sitting in cohomological degree $0$.
\end{enumerate}
Then there is a canonical bivariant class
$$
P'_p(E_2),\text{ resp. }c'_p(E_2) \in A^p(X_2 \to X)
$$
characterized by the property
$$
P'_p(E_2) \cap i_{2, *} \alpha_2 = P_p(E_2) \cap \alpha_2
\quad\text{and}\quad
P'_p(E_2) \cap i_{1, *} \alpha_1 = 0,
$$
respectively
$$
c'_p(E_2) \cap i_{2, *} \alpha_2 = c_p(E_2) \cap \alpha_2
\quad\text{and}\quad
c'_p(E_2) \cap i_{1, *} \alpha_1 = 0
$$
for $\alpha_i \in \CH_k(X_i)$ and similarly after any base change
$X' \to X$ locally of finite type.
\end{lemma}

\begin{proof}
We are going to use the material of Section \ref{section-pre-derived}
without further mention.

\medskip\noindent
Assume $E_2|_{X_1 \cap X_2}$ is zero.
Consider a morphism of schemes $X' \to X$
which is locally of finite type and denote $i'_j : X'_j \to X'$ the
base change of $i_j$. By Lemma \ref{lemma-exact-sequence-closed-chow}
we can write any element $\alpha' \in \CH_k(X')$ as
$i'_{1, *}\alpha'_1 + i'_{2, *}\alpha'_2$ where
$\alpha'_2 \in \CH_k(X'_2)$
is well defined up to an element in the image of pushforward
by $X'_1 \cap X'_2 \to X'_2$. Then we can set
$P'_p(E_2) \cap \alpha' = P_p(E_2) \cap \alpha'_2 \in \CH_{k - p}(X'_2)$. This
is well defined by our assumption that $E_2$ restricts
to zero on $X_1 \cap X_2$.

\medskip\noindent
If $E_2|_{X_1 \cap X_2}$ is isomorphic to a finite locally free
$\mathcal{O}_{X_1 \cap X_2}$-module of rank $< p$ sitting in
cohomological degree $0$, then $c_p(E_2|_{X_1 \cap X_2}) = 0$
by rank considerations and we can argue in exactly the same manner.
\end{proof}

\end{document}
